\documentclass[12pt, a4paper]{article}
\usepackage[T1]{fontenc}
\usepackage{lmodern}
\usepackage{amsmath, amsthm, amssymb, mathtools}
\usepackage[margin=1in]{geometry}
\usepackage{xr-hyper}
\usepackage[colorlinks=true,linkcolor=blue,citecolor=blue]{hyperref}
\usepackage{cleveref}
\usepackage{graphicx, subcaption}
\usepackage{natbib}
\usepackage{booktabs}
\usepackage{enumerate}

% Code references: scaled monospace with line-breaking
\newcommand{\code}[1]{\texttt{\small #1}}

\newtheorem{theorem}{Theorem}[section]
\newtheorem{lemma}[theorem]{Lemma}
\newtheorem{proposition}[theorem]{Proposition}
\newtheorem{corollary}[theorem]{Corollary}
\newtheorem{definition}{Definition}[section]
\newtheorem{remark}{Remark}[section]
\newtheorem{conjecture}[theorem]{Conjecture}
\newtheorem{derivation}[theorem]{Derivation}
\newtheorem{observation}[theorem]{Observation}
\numberwithin{equation}{section}

\newcommand{\dd}{\mathrm{d}}
\newcommand{\seriestitle}{Why Rotating Fluids Make Polygons}


\externaldocument[I-]{../paper-1-mathematics/main}
\externaldocument[II-]{../paper-2-physics/main}
\externaldocument[III-]{../paper-3-gravity/main}
\externaldocument[A-]{../paper-A-appendices/main}

\title{\seriestitle\\[6pt]\large Paper IV: Planetary and Astrophysical Applications}
\author{Gordon Speagle}
\date{}

\begin{document}
\maketitle

\begin{abstract}
Applying the constraint-intersection principle from Papers~I--II, we predict the polygon wavenumber for Saturn ($N=6$, Rossby-limited), Jupiter north ($N=8$, Thomson-limited), Jupiter south ($N=5$, packing-limited), and Neptune ($N \in \{3,4,5\}$, packing-dominated). Saturn's $10\%$ oblateness is not a correction but the mechanism: it flips the Thomson eigenvalue $\lambda_3$ from $-0.12$ (unstable on a mean sphere) to $+0.01$ (marginally stable), with a critical latitude $\varphi_{\mathrm{crit}} \approx 72^\circ$N below which the hexagon cannot exist. We derive a first-principles prediction chain from atmospheric parameters and resolve the Onsager--Arnold tension via a two-timescale argument.
\end{abstract}

\label{part:planets}
%% ====================================================================

%% ====================================================================
\section{The three-scale selection mechanism}
\label{sec:three-scale}
\label{sec:two-scale}
\label{sec:extensions}

Three scales cooperate:
\emph{macro} (Onsager inverse cascade concentrates vorticity into
$N$ clusters---a physical hypothesis, not derived),
\emph{meso} (the constraint-intersection principle selects which $N$),
and \emph{micro} (Arnold/Thomson stability arranges the clusters as
a regular polygon).
The Onsager hypothesis is supported by the \citet{Siegelman2022}
simulations ($N = 8$ north, $N = 5$ south in 2D turbulence)
but remains empirically untested for Jupiter's actual 3D dynamics.


%% ====================================================================
\section{Wavenumber selection: the constraint-intersection principle}
\label{sec:constraint-intersection}

\begin{proposition}[Constraint-intersection selection]\label{prop:complementary}
The observed polygon wavenumber $N$ is determined by the intersection of
three constraints.  Two are fully derived from the 2D Laplacian:
\begin{enumerate}
  \item \textbf{Rossby stationarity} (derived; for jet meanders):
    the QGPV eigenvalue problem
    $(\nabla^2 + \beta\,\partial_y)\psi = 0$ selects the stationary
    wavenumber $n^*$.
  \item \textbf{Thomson stability} (derived; for discrete cyclone rings):
    $\kappa_0/\kappa \ge \kappa_{\mathrm{crit}}(N)$ gives an upper
    bound $N \le N_{\max}(\kappa_0/\kappa)$; the energy monotonicity
    $H_{N+1} > H_N$ (Proposition~\ref{prop:onsager-n-select}) then
    selects $N = N_{\max}$.
\end{enumerate}
The third has a derived skeleton plus one empirical input:
\begin{enumerate}\setcounter{enumi}{2}
  \item \textbf{Geometric packing} (derived skeleton + empirical factor):
    domain admissibility of $G(\nabla^2)$ on the punctured domain
    gives $\sin(\pi/N) \ge \varepsilon/R$ where $\varepsilon$ is
    the vortex core radius (derived).  The effective exclusion radius
    $r_{\mathrm{excl}} \approx 1.3\,\varepsilon$ includes a
    $\beta$-drift buffer zone from nonlinear vortex dynamics
    \citep{Gavriel2021}; the factor $1.3$ is observational, not
    derivable from the linear operator.
\end{enumerate}
The selected wavenumber is
\begin{equation}\label{eq:constraint-intersection}
  N_{\mathrm{selected}} = \max\bigl\{N :
    N \le N_{\mathrm{Thomson}},\;
    \sin(\pi/N) \ge r_{\mathrm{excl}}/R,\;
    n\;\text{Rossby-stationary}\bigr\},
\end{equation}
where inapplicable constraints are dropped (e.g.\ Rossby stationarity
does not apply to discrete cyclone rings, and packing does not constrain
jet meanders).
\end{proposition}

\begin{remark}[Constraint-intersection principle]
\label{rem:constraint-intersection}
Each planetary system activates a different subset of the three
constraints in~\eqref{eq:constraint-intersection}:
\begin{center}
\begin{tabular}{lcccl}
\toprule
System & $N_{\mathrm{Thomson}}$ & $N_{\mathrm{Rossby}}$ & $N_{\mathrm{pack}}$ & Binding \\
\midrule
Saturn (76${}^\circ$N)  & $7$ & $6$ & ---  & Rossby \\
Jupiter north           & $8$ & --- & $>8$ & Thomson \\
Jupiter south           & $8$ & --- & $5$  & Packing \\
\bottomrule
\end{tabular}
\end{center}
The framework is unified: all three constraints derive from the 2D
Laplacian~$\nabla^2$ and its modifications.  The packing constraint
is \emph{not} an external geometric condition: it is the domain
admissibility condition for the Green's function
$G_\varepsilon$ of~$\nabla^2$ on the punctured domain
$\Omega_N(\varepsilon) = \mathbb{R}^2 \setminus
\bigcup_{k=1}^N B(z_k, \varepsilon)$, where $\varepsilon$ is the
vortex core radius.  This Green's function exists and defines a
well-posed BVP if and only if the excluded discs are pairwise disjoint:
\[
  |z_j - z_k| \ge 2\varepsilon
  \;\Longleftrightarrow\;
  \sin(\pi/N) \ge \varepsilon/R.
\]
The domain-admissibility condition establishes
$r_{\mathrm{excl}} \ge \varepsilon$ (the core radius) from the
Laplacian alone.  The observed enlargement
$r_{\mathrm{excl}} \approx 1.3\,\varepsilon$ at Jupiter's south pole
arises from the nonlinear $\beta$-drift balance of finite-amplitude
vortices on a rotating sphere \citep{Gavriel2021}; this factor is
\emph{not} derived from the linear operator $\nabla^2 + \beta\,\partial_y$
but is an empirical input from the full vortex dynamics.

The two Laplacian-derived facets and their observational supplement are:
\begin{enumerate}
  \item \emph{Spectral positivity} of $G(\nabla^2)$: Havelock
    eigenvalue $\lambda_m > 0$ $\to$ Thomson bound (derived);
  \item \emph{Domain admissibility} of $G_\varepsilon(\nabla^2)$:
    disc non-overlap $\sin(\pi/N) \ge \varepsilon/R$ $\to$ packing
    bound with $r_{\mathrm{excl}} = \varepsilon$ (derived);
    the enlargement $r_{\mathrm{excl}} \approx 1.3\,\varepsilon$
    at Jupiter south is observational;
  \item \emph{Rossby stationarity} of $(\nabla^2 + \beta\,\partial_y)\psi = 0$:
    stationary mode $n^*$ $\to$ wave quantization (derived).
\end{enumerate}
The finite-core regularisation $G_\varepsilon$ does \emph{not} create
a blob-stable window for $N = 8$: the corrected constrained Hessian
(with the full $\partial^2(-\tfrac12\ln s)/\partial x_j^2
= (2\,\Delta x^2 - s)/s^2$, $s = d^2 + 2\varepsilon^2$, not
the $O(\varepsilon^2)$ truncation) remains negative throughout
the admissible domain $\varepsilon/R < \sin(\pi/8)$.
The $N = 8$ selection therefore requires the central vortex
($\kappa_0/\kappa \ge \tfrac{1}{2}$); blob regularisation alone
does not provide an alternative stabilisation mechanism.
\end{remark}

For Saturn, $n^*$ is determined by the Rayleigh--Kuo condition for
stationary waves on a jet.  All quantities are evaluated in the frame
rotating with Saturn's System~III period---the frame co-rotating with
the deep interior.  In this frame the background wind at the hexagonal
jet is $U_{\max} \approx 120$\,m/s eastward.  At the center of a
Gaussian jet, the effective PV gradient is
$\beta_{\mathrm{eff}} = \beta + U_{\max}/\sigma_{\mathrm{jet}}^2$,
where the jet curvature term dominates bare $\beta$ by a factor
$\sim\!14$ for $\sigma_{\mathrm{jet}} = 2.5 \times 10^6$\,m
\citep{Pedlosky1987, Vallis2017}.
Two estimates of $n^*$ are available:
\begin{itemize}
  \item \emph{Simple scaling} (independent of $\sigma_{\mathrm{jet}}$):
    $n^* = R_{\mathrm{hex}}\sqrt{\beta/U_{\max}} = 5.6$, which
    rounds to $n = 6$ for any $\sigma_{\mathrm{jet}}$.
  \item \emph{Leading-order with $\beta_{\mathrm{eff}}$}:
    $n^* = R\cos\varphi\,\sqrt{\beta_{\mathrm{eff}}/U_{\max}}$,
    which gives $5.4$ at $\sigma_{\mathrm{jet}} = 2.5 \times 10^6$\,m
    but is sensitive to the jet width ($n^* = 6.4$ at
    $\sigma_{\mathrm{jet}} = 2.1 \times 10^6$\,m).
\end{itemize}
The full numerical Rayleigh--Kuo eigenvalue computation yields
$n^* = 6$ \citep{SanchezLavega2014}; the simple scaling reproduces
this robustly, while the $\beta_{\mathrm{eff}}$ formula is sensitive
to $\sigma_{\mathrm{jet}}$ within its $\sim\!20\%$ observational
uncertainty.  We adopt
$\sigma_{\mathrm{jet}} = 2.5 \times 10^6$\,m throughout, consistent
with the Cassini jet profile measurements; all conclusions are
verified to hold over the range
$\sigma_{\mathrm{jet}} \in [2.0, 3.0] \times 10^6$\,m.
Thomson allows $N \le 7$; both criteria are satisfied at $N = 6$.

For Jupiter's north pole ($N = 8$): the cyclones are discrete vortex
clusters, not a jet meander; Rossby stationarity gives
$n^* \approx 0.3$ (not physically relevant).  \emph{Caveat:} the 2D
point vortex model requires the cyclone roots to extend through a
barotropic column much deeper than the ring radius.  Juno gravity
data constrains Jupiter's differential rotation to
$\sim\!3{,}000$\,km depth \citep{Kaspi2018}, roughly $4\%$ of
$R_J$---making Jupiter the least observationally supported case for
the 2D assumption among our applications.  The model is nevertheless
applicable if the polar cyclones are confined to a deep weather layer
($H_{\mathrm{eff}} \sim 3{,}000$\,km, $R_d \gg R$) or to a shallow
tropospheric layer with reduced gravity; both scenarios are consistent
with the Juno data (see Section~Paper~II, \S13 for
quantitative limits).

Thomson stability provides an \emph{upper bound}:
$\kappa_{\mathrm{crit}}(9) = 1 > \kappa_0/\kappa$,
so $N = 9$ is unstable; $\kappa_{\mathrm{crit}}(8) = \tfrac{1}{2}
\le \kappa_0/\kappa$, so $N = 8$ is stable \citep{Adriani2018}.
But $N = 7$ is also stable ($\kappa_{\mathrm{crit}}(7) = 0$), so
stability alone does not select $N = 8$ over $N = 7$.
The selection requires one additional ingredient: a reason for the
system to prefer higher-energy configurations among the stable set.

\begin{proposition}[Energy monotonicity and maximum-$N$ selection]\label{prop:onsager-n-select}
Let $\Gamma = N\kappa$ be the total ring circulation, $L = N\kappa R^2$
the angular impulse (so $R$ is fixed when $\Gamma$ and $L$ are fixed),
and $r = \kappa_0/\kappa$ held fixed as $N$ varies (both the ring and
central circulations scale proportionally).  Denote by
$r_{\mathrm{core}}$ the vortex core radius---the minimum separation at
which the point-vortex model is valid; equivalently, $r_{\mathrm{core}}$
is the unit of length in which the logarithm is taken, so that
$\ln R > 0$ iff $R > r_{\mathrm{core}}$.  Then
\begin{equation}\label{eq:H-monotone}
  H_{N+1} - H_N = \Gamma^2\!\left[\frac{2r-1}{2N(N+1)}\ln R
    + \frac12\!\left(\frac{\ln N}{N} - \frac{\ln(N+1)}{N+1}\right)\right].
\end{equation}
For $r \ge \tfrac{1}{2}$ and $R > r_{\mathrm{core}}$, both terms
are non-negative (first: $2r - 1 \ge 0$ and $\ln R > 0$; second:
$\ln N/N$ is strictly decreasing for $N > e \approx 2.7$), so
$H_{N+1} > H_N$ for all $N \ge 3$.  Therefore the $N$-gon ring
energy is strictly monotone in $N$: the maximum-$H$ configuration
within the Thomson-stable set is the ring with the largest $N$:
\[
  N_{\max}(\kappa_0/\kappa)
    \equiv \max\bigl\{N' : \kappa_{\mathrm{crit}}(N') \le \kappa_0/\kappa\bigr\}.
\]
Any dynamics that preferentially populates higher-energy stable
configurations---whether Onsager negative-temperature equilibrium
($\beta < 0$, where the Boltzmann weight $e^{-\beta H}$ increases
with $H$), sequential vortex injection (adding cyclones until
Thomson instability prevents $N{+}1$; this requires fixed~$r$,
see Remark~\ref{rem:fixed-r}), or merger of smaller vortex
clusters---selects $N = N_{\max}$.  The mathematical content of the
proposition is the monotonicity $H_{N+1} > H_N$; the physical
mechanism by which the system reaches $N_{\max}$ is a separate
question (see the discussion below and Section~\ref{sec:extensions}).

The stability threshold $\kappa_{\mathrm{crit}}(8) = \tfrac{1}{2}$ is
exactly the value at which $H_{N+1} > H_N$ becomes unconditional (i.e.\
$R$-independent): at $r = \tfrac{1}{2}$, term~1 vanishes identically
and term~2 remains strictly positive.
\end{proposition}
\begin{proof}
Use the Hamiltonian
$H = -\kappa^2\!\sum_{j<k}\ln|z_j - z_k| - \kappa_0\kappa\!\sum_j\ln|z_j|$
with $\kappa = \Gamma/N$ and $\kappa_0 = r\kappa = r\Gamma/N$.

\emph{Step 1: Product formula.}
For the regular $N$-gon $z_j = Re^{2\pi ij/N}$, the chord length
between vortices $j$ and $k$ is $|z_j - z_k| = 2R\sin(\pi|j-k|/N)$.
By $\mathbb{Z}_N$ symmetry the pair sum equals $N/2$ times the sum from
one vertex:
\[
  \sum_{j<k}\ln|z_j - z_k|
  = \frac{N}{2}\sum_{s=1}^{N-1}\ln\!\bigl(2R\sin(\pi s/N)\bigr)
  = \frac{N}{2}\bigl[\ln N + (N-1)\ln R\bigr].
\]
Here we used $\sum_{s=1}^{N-1}\ln(2\sin(\pi s/N)) = \ln N$, which
follows from evaluating the polynomial identity $(z^N-1)/(z-1) =
\prod_{k=1}^{N-1}(z - e^{2\pi ik/N})$ at $z = 1$: the product of
moduli is $\prod_{s=1}^{N-1}|1 - e^{2\pi is/N}|
= \prod_{s=1}^{N-1}2\sin(\pi s/N) = N$.

\emph{Step 2: Central vortex term.}
Since $|z_j| = R$ for all $j$: $\sum_j\ln|z_j| = N\ln R$.

\emph{Step 3: Closed form for $H_N$.}
Substituting Steps~1--2 with $\kappa = \Gamma/N$, $\kappa_0 = r\Gamma/N$:
\begin{equation}\label{eq:HN-explicit}
  H_N = -\frac{\Gamma^2}{2N}\ln N
        - \Gamma^2\!\left(\frac{N-1}{2N} + \frac{r}{N}\right)\ln R.
\end{equation}

\emph{Step 4: Difference $H_{N+1} - H_N$.}
Writing $H_{N+1}$ from \eqref{eq:HN-explicit} with $N\to N+1$ and
subtracting:
\begin{align*}
  H_{N+1} - H_N
  &= -\frac{\Gamma^2}{2}\!\left[\frac{\ln(N{+}1)}{N{+}1}
     - \frac{\ln N}{N}\right]
     - \Gamma^2\ln R\!\left[\frac{N/2+r}{N{+}1}
     - \frac{(N{-}1)/2+r}{N}\right].
\end{align*}
The coefficient of $\ln R$ simplifies by partial fractions:
\[
  \frac{N/2+r}{N{+}1} - \frac{(N{-}1)/2+r}{N}
  = \frac{N^2 - (N{-}1)(N{+}1)}{2N(N{+}1)} - \frac{r}{N(N{+}1)}
  = \frac{1-2r}{2N(N{+}1)},
\]
where we used $N^2-(N-1)(N+1) = 1$.  Negating and rearranging
gives~\eqref{eq:H-monotone}.

\emph{Step 5: Sign of each term.}
\emph{Term~2}: $f(x) = \ln x/x$ satisfies $f'(x) = (1-\ln x)/x^2 < 0$
for $x > e$, so $\ln N/N > \ln(N+1)/(N+1)$ for all $N \ge 3$.

\emph{Term~1}: $2r-1 \ge 0$ for $r \ge \tfrac{1}{2}$; and
$\ln R > 0$ for $R > r_{\mathrm{core}}$.  Hence both terms are
non-negative, and term~2 is strictly positive, giving
$H_{N+1} > H_N$.
\end{proof}

\begin{remark}[Physical status of the fixed-$r$ assumption]
\label{rem:fixed-r}
The assumption $r = \kappa_0/\kappa$ fixed as $N$ varies means that
as more cyclones appear in the ring ($N$ increases), each individual
ring vortex becomes weaker ($\kappa = \Gamma/N$) \emph{and the central
vortex weakens proportionally} ($\kappa_0 = r\Gamma/N$).  This is
non-trivial: the central polar cyclone is a distinct dynamical object
from the ring cyclones and need not respond to changes in $N$.

The assumption is justified when all vortices --- ring and central ---
are drawn from the same turbulent condensation process, so their
strengths are set by a common large-scale energy budget rather than
by independent forcings.  In the Onsager picture, the negative-temperature
state is reached by the inverse cascade from the same vorticity
reservoir, giving a fixed partition of total circulation between center
and ring.  If instead $\kappa_0$ is held fixed while $\kappa = \Gamma/N$
varies (independent central vortex), then $r = N\kappa_0/\Gamma$ grows
with $N$ and the monotonicity of $H_{N+1} > H_N$ strengthens further,
so the maximum-$N$ selection remains qualitatively intact.  The
marginal case is $\kappa_0 = 0$ (no central vortex): then
$r = 0 < \tfrac12$ and Term~1 in~\eqref{eq:H-monotone} is
\emph{negative}, so $H_{N+1} > H_N$ requires $R$ large enough that
Term~2 dominates.  The proposition's $r \ge \tfrac{1}{2}$ condition
is therefore specific to the regime in which the central vortex is
sufficiently strong relative to the ring vortices.

The selection of $N = 8$ over $N = 7$ is sensitive to which assumption
is made.  Under fixed~$r$: $H_{N+1} > H_N$ for all $N \ge 3$
(Proposition~\ref{prop:onsager-n-select}), so Onsager selects the
Thomson-maximum $N_{\max}$.  Under fixed~$\kappa_0$: $r$ grows with
$N$, which strengthens the monotonicity and gives the same $N_{\max}$.
Under the intermediate assumption $\kappa_0 = \mathrm{const}$ while
$\Gamma$ also varies: the result depends on how $\Gamma(N)$ scales,
which is not observationally constrained.  Thus the $N = 8$
selection is robust across the two limiting assumptions (fixed~$r$
and fixed~$\kappa_0$) but relies on the Onsager framework providing
the link between $N$, $\Gamma$, and $\kappa_0$---a physical
hypothesis, not a mathematical consequence.

Observationally, Jupiter's central cyclone differs from the ring
cyclones: it has a smaller diameter
($\sim\!4{,}000$\,km vs.\ $4{,}000$--$7{,}000$\,km;
\citealt{Adriani2018}) and appears more axisymmetric.  Whether it
formed from the same inverse-cascade process as the ring cyclones,
or was generated by a distinct polar convective mechanism, is not
resolved by current Juno data.
\end{remark}

\noindent
For Jupiter north, $\kappa_0/\kappa \approx 1$ with roughly 30\%
observational uncertainty (Juno; \citealt{Adriani2018}), giving a
conservative range $\kappa_0/\kappa \in [0.7, 1.3]$.  Since
$\kappa_{\mathrm{crit}}(8) = \tfrac{1}{2}$ exactly, this entire
range satisfies the stability condition for $N = 8$.  However,
$\kappa_{\mathrm{crit}}(9) = 1$, which lies \emph{inside} the
uncertainty interval: if $\kappa_0/\kappa \ge 1$ (the upper half
of the range), then $N = 9$ is also Thomson-stable and the Onsager
argument would select $N_{\max} = 9$ instead.  The conclusion
$N_{\max} = 8$ therefore relies on $\kappa_0/\kappa < 1$, which is
consistent with the central estimate but not guaranteed by the
current observational precision.  More precisely:
$\kappa_0/\kappa \le 1$ gives $N_{\max} = 8$ (observed);
$\kappa_0/\kappa > 1$ gives $N_{\max} = 9$ (not observed).  The
observed $N = 8$ is thus evidence that $\kappa_0/\kappa \lesssim 1$,
providing a theory-derived constraint on the central vortex
strength that could be tested with improved Juno wind
measurements.

\label{sec:n-selection}%
For Jupiter's south pole ($N = 5$): $\kappa_{\mathrm{crit}}(N) = 0$
for all $N \le 7$, so the Thomson bound alone allows up to $N = 7$
regardless of the central vortex.  The observed $N = 5$ is selected
by the geometric packing constraint
(Proposition~\ref{prop:complementary}, item~3).
South-pole cyclones are individually $\approx 2\times$ larger in
linear size than north-pole cyclones ($r_{\mathrm{core}} \approx
3{,}500$\,km vs.\ $2{,}500$\,km; \citealt{Adriani2018}), and each
cyclone carries a $\beta$-drift exclusion zone that extends $\approx
30\%$ beyond the core radius \citep{Gavriel2021}, giving
$r_{\mathrm{excl}} \approx 1.3\,r_{\mathrm{core}} \approx 4{,}550$\,km.
At the south-pole ring radius $R \approx 8{,}000$\,km
($\sim\!7^{\circ}$ from the pole):
\[
  \sin(\pi/5) = 0.588 \ge \frac{r_{\mathrm{excl}}}{R}
  = \frac{4{,}550}{8{,}000} = 0.569,
  \qquad
  \sin(\pi/6) = 0.500 < 0.569.
\]
Thus $N_{\mathrm{pack}} = 5$, which is tighter than the Thomson bound
$N_{\mathrm{Thomson}} = 8$: packing, not stability, is the binding
constraint at Jupiter's south pole.  The $\approx 30\%$ excess of
$r_{\mathrm{excl}}$ over the bare cyclone radius is physically the
$\beta$-drift buffer zone identified by \citet{Gavriel2021}---cyclone
repulsion on a rotating sphere creates an effective ``personal space''
larger than the vortex itself.


%% ====================================================================
\section{Saturn's hexagon}
\label{sec:saturn}

\subsection{Eigenvalue branch disconnection (Saturn persistence)}

The preceding subsection selected Saturn's wavenumber $n = 6$ via
Rossby stationarity.  A natural follow-up: can smooth perturbations
of the jet destroy the stationary wave, converting the hexagonal
meander into an evanescent (decaying) mode?  The answer is no, by
a WKB branch-disconnection argument.

The Rayleigh--Kuo equation for stationary ($c = 0$) mode~$n$ on a
zonal jet $U(y)$ reads $\hat\psi'' + V(\rho)\hat\psi = 0$, where
$V = (\beta - U'')/U - n^2/R^2$ is the effective potential in
log-polar coordinates.  At the jet center ($y = 0$) of a Gaussian jet:
\[
  V(0) = \frac{\beta + U_{\max}/\sigma^2}{U_{\max}} - \frac{n^2}{R^2}
  = \frac{\beta_{\mathrm{eff}}}{U_{\max}} - \frac{n^2}{R^2}.
\]
For Saturn ($U_{\max} = 120\;\mathrm{m/s}$, $\sigma = 2.5\times
10^6\;\mathrm{m}$, $\beta = 1.46\times 10^{-12}\;\mathrm{m^{-1}
s^{-1}}$, $n = 6$, $R = R_{\mathrm{S}}\cos 76^\circ = 1.31\times
10^7\;\mathrm{m}$):
$V(0) = 1.72\times 10^{-13} - 2.10\times 10^{-13} =
-3.8\times 10^{-14}\;\mathrm{m}^{-2}$ (negative).
The sign of $V(0)$ depends on $\sigma_{\mathrm{jet}}$:
$V(0) < 0$ for $\sigma_{\mathrm{jet}} \gtrsim 2.3 \times 10^6$\,m,
$V(0) > 0$ for narrower jets.  At the adopted
$\sigma_{\mathrm{jet}} = 2.5 \times 10^6$\,m, $V(0) < 0$ with
$|V(0)|/V_{\mathrm{threshold}} = 1.6$ (a $60\%$ margin above the
sign change); this margin grows to $3.4\times$ at
$\sigma_{\mathrm{jet}} = 3.0 \times 10^6$\,m.  The Cassini
jet profile measurements constrain
$\sigma_{\mathrm{jet}} \in [2.3, 3.0] \times 10^6$\,m, placing
$V(0)$ robustly in the negative regime.

Since $V < 0$, the WKB exponent $s^2 = V < 0$ forces oscillatory
solutions with $s = i\alpha_\rho$ ($\alpha_\rho \ne 0$).  Writing
$s = \sigma + i\alpha_\rho$:
$\operatorname{Im}(s^2) = 2\sigma\alpha_\rho = 0$.
When $\alpha_\rho \ne 0$ (oscillatory region): $\sigma = 0$.
The elliptic ($\sigma = 0$) and hyperbolic ($\alpha_\rho = 0$)
branches are disconnected: smooth jet perturbations cannot move the
solution from the oscillatory branch to the evanescent branch at
leading WKB order.  This is a structural stability result for the
Saturn selection mechanism: the $n = 6$ stationary wave exists
because $V(0) < 0$, and no continuous deformation of the jet profile
can remove it without first passing through $V(0) = 0$ (which
requires changing the sign of $\beta_{\mathrm{eff}}/U_{\max} -
n^2/R^2$, a finite perturbation).  The hexagon is therefore
structurally stable at leading WKB order: it cannot be removed by
infinitesimal jet perturbations.

\subsection{The oblateness correction (Saturn's shape selects N=6)}
\label{sec:oblateness}

The analysis above treats Saturn as a sphere.  Saturn is an oblate
spheroid with flattening $f = (a{-}b)/a = 0.098$
($a = 60{,}268$\,km equatorial, $b = 54{,}364$\,km polar).
The oblateness changes the \emph{sign} of the Thomson stability
eigenvalue for $N = 6$---it is not a small correction but the
mechanism that determines whether the hexagon can exist.

\paragraph{Curvature on the oblate spheroid.}
The Gaussian curvature at geodetic latitude $\varphi$ is
$K(\varphi) = 1/(M(\varphi)\,N(\varphi))$ where
$M = a(1{-}e^2)/(1{-}e^2\sin^2\!\varphi)^{3/2}$ and
$N = a/(1{-}e^2\sin^2\!\varphi)^{1/2}$
($e^2 = 1 - b^2/a^2$).
On an oblate body, $K$ is \emph{largest} at the equator and
\emph{smallest} at the pole---the pole is flatter.
For Saturn: $K_{\mathrm{eq}}/K_{\mathrm{pole}} = a^4/b^4 = 1.51$.
Since positive curvature destabilises (Paper~I,
Proposition~\ref{I-prop:curvature-dichotomy}), the pole is
the \emph{most stable} point on the spheroid---closer to the
flat-plane limit where $N_{\mathrm{crit}} = 7$.

\paragraph{The critical result.}
Substituting the local curvature $K(\varphi)$ into the spherical
$C_1$ formula (a local-curvature approximation valid for
$K R^2 \ll 1$):
\begin{center}
\begin{tabular}{lcc}
\toprule
Model & $\lambda_3(N{=}6)$ & Stable? \\
\midrule
Mean sphere ($R_{\mathrm{mean}} = 58{,}300$\,km) &
  $-0.12$ & No \\
Oblate spheroid at $78^\circ$N &
  $+0.01$ & Marginally yes \\
\bottomrule
\end{tabular}
\end{center}
The oblateness flips $\lambda_3$ from $-0.12$ to $+0.01$.
On a mean sphere, the hexagon is Thomson-unstable and
\emph{should not exist}; on the oblate spheroid, it is
marginally stable.

\paragraph{Critical latitude.}
The latitude $\varphi_{\mathrm{crit}}$ at which $\lambda_3 = 0$
is determined by $C_1(K(\varphi_{\mathrm{crit}})\cdot R^2) =
f(3,6) = 9/2$.  For $R_{\mathrm{ring}} = 15{,}000$\,km:
$\varphi_{\mathrm{crit}} \approx 72^\circ\mathrm{N}$.
Saturn's hexagon at $78^\circ$N is $6^\circ$ above this
critical latitude---comfortably in the stable region but
not far from the boundary.

\begin{proposition}[Oblateness-stabilised hexagon]
\label{prop:oblateness}
On an oblate spheroid with flattening $f$, the binding
eigenvalue $\lambda_3(N{=}6)$ at the pole exceeds its mean-sphere
value by $\delta\lambda_3 \approx (N{-}1)\cdot 2\xi\, f\,(2{-}f)
/(1{+}\xi)^2$ at leading order in $f$, where $\xi = K_{\mathrm{mean}}
R^2$.  For Saturn ($f = 0.098$, $\xi \approx 0.07$):
$\delta\lambda_3 \approx 0.13$, which reverses the sign of
$\lambda_3$ from $-0.12$ to $+0.01$.
\end{proposition}

\paragraph{Mode coupling at non-polar latitudes.}
At $78^\circ$N, the Gaussian curvature varies by
$\delta K/K \approx 4.2\%$ across the hexagon ring (from the
curvature gradient $\dd K/\dd\varphi$).  This $\cos(2\theta)$
variation couples Fourier modes $m$ and $m \pm 2$, with coupling
amplitude $\sim (\delta K/K)\cdot C_1 \approx 0.19$.  Since the
stability margin is only $|\lambda_3| = 0.01$, the mode coupling
\emph{exceeds} the margin by a factor of~$19$.

By the topological protection of the Morse index
(Paper~I, Remark~\ref{I-rmk:topological-protection}), this
mode coupling cannot change the \emph{net} spectral flow for
$N = 6$ (which has all eigenvalues strictly positive at the pole).
Individual eigenvalue trajectories are perturbed, but the
net count of instabilities is zero---the hexagon remains stable
in the topological sense, even though the perturbative stability
margin is formally violated.

\paragraph{Neptune and Uranus: oblateness is negligible.}
Neptune ($f = 0.017$) and Uranus ($f = 0.023$) are far less
oblate.  The polar curvature correction $\delta K/K \approx
-4.5\%$ (Neptune) and $-6.8\%$ (Uranus) shifts $\lambda_3$
by $\lesssim 0.01$---negligible compared to the Thomson
eigenvalue at these planets' ring radii.  The Neptune prediction
$N \in \{3, 4, 5\}$ (Section~\ref{sec:observations}) is robust
to spheroidal corrections
(\texttt{extensions/oblate\_spheroid.py}).


%% ====================================================================
\section{Jupiter's polar polygons}
\label{sec:jupiter}

\subsection{North pole: $N = 8$ from Thomson + central vortex}

The north-pole octagon ($N = 8$) is Thomson-limited.
On the flat plane, $N_{\mathrm{crit}} = 7$; the octagon's
critical eigenvalue is $\lambda_4 = 7 - 8 = -1 < 0$ (unstable).
A central vortex of strength $\kappa_0$ shifts this by
$+2\kappa_0$; the critical ratio is
$\kappa_{\mathrm{crit}}(8) = |\lambda_4|/2 = 1/2$.
For $N = 9$: $\kappa_{\mathrm{crit}}(9) = 1$; for $N = 10$:
$\kappa_{\mathrm{crit}}(10) = 7/4$.
Juno imagery shows a central cyclone of roughly equal strength
to the ring cyclones ($\kappa_0/\kappa \approx 1$),
which stabilises $N = 8$ but not $N = 9$.
The energy monotonicity $H_{N+1} > H_N$
(Proposition~\ref{prop:onsager-n-select})
then selects $N = 8$ as the Onsager maximum.

\smallskip\noindent
The $K_0$ interaction constrains the weather-layer depth.
At the Rossby deformation radius $R_d \approx 2{,}500$\,km
(matching the observed cyclone radius at
$H_{\mathrm{eff}} \approx 3{,}000$\,km,
$g' \approx 0.3$\,m\,s$^{-2}$), the ring-to-deformation ratio
$R_{\mathrm{ring}}/R_d \approx 8{,}500/2{,}500 \approx 3.4$.
The $K_0$ stability analysis gives:
\begin{center}
\small
\begin{tabular}{lcl}
\toprule
$R/R_d$ & $N_{\mathrm{crit}}$ & Regime \\
\midrule
$0$ & $7$ & Barotropic limit \\
$0.8$ & $8$ & Deep weather layer \\
$1.0$ & $8$ & $\leftarrow$ requires $R_d \gtrsim R_{\mathrm{ring}}$ \\
$1.5$ & $7$ & \\
$2.0$ & $6$ & Shallow weather layer \\
\bottomrule
\end{tabular}
\end{center}
\noindent
$N = 8$ requires $R/R_d \lesssim 1$, i.e.\
$R_d \gtrsim 8{,}500$\,km.
This constrains the effective weather-layer depth:
$H_{\mathrm{eff}} = f^2 R_d^2/g' \gtrsim 3{,}000$\,km,
a testable prediction for Juno gravity measurements.

\subsection{South pole: $N = 5$ from geometric packing}

The south-pole pentagon ($N = 5$) is packing-limited.
With $R_{\mathrm{ring}} \approx 8{,}500$\,km,
cyclone radius $r_{\mathrm{cyclone}} \approx 3{,}500$\,km, and
$\beta$-drift buffer $\approx 1.3$ \citep{Gavriel2021}:
the effective exclusion radius
$r_{\mathrm{excl}} \approx 4{,}550$\,km gives
$r_{\mathrm{excl}}/R_{\mathrm{ring}} \approx 0.535$.
Since $\sin(\pi/6) = 0.500 < 0.535$: a hexagonal ring does
\emph{not} fit---the cyclones would overlap.
Since $\sin(\pi/5) = 0.588 > 0.535$: a pentagonal ring fits.
Therefore $N_{\mathrm{pack}} = 5$.

The selection is \emph{tight}: it requires
$R_{\mathrm{ring}} < 2\,r_{\mathrm{excl}} = 9{,}100$\,km
for $N = 6$ to be excluded.  The observed
$R_{\mathrm{ring}} \approx 8{,}500$\,km satisfies this with
$7\%$ margin.  A testable prediction: if the south-pole ring
radius were $> 9{,}100$\,km, the framework would predict $N = 6$
instead of~$5$.

At the north pole: $r_{\mathrm{excl}}/R \approx 0.38$,
so $N_{\mathrm{pack}} \ge 8$---packing is not the binding
constraint.

\subsection{The north--south asymmetry in one parameter}

The entire $N = 8$ vs.\ $N = 5$ asymmetry reduces to a single
dimensionless ratio:
\[
  \chi = \frac{R_d^{(\mathrm{south})}}{R_d^{(\mathrm{north})}}
  \approx \frac{r_{\mathrm{cyclone}}^{(\mathrm{south})}}
              {r_{\mathrm{cyclone}}^{(\mathrm{north})}}
  \approx \frac{3{,}500}{2{,}500} = 1.40.
\]
The constraint intersection requires $\chi \gtrsim 1.30$
for the south-pole packing bound to tighten from $N = 6$ to
$N = 5$.  The observed $\chi = 1.40$ matches to $8\%$.
Deriving $\chi$ from atmospheric physics (insolation asymmetry,
moist convection) remains open: radiative scaling gives
$\chi \sim 1.07$ (insufficient), suggesting that deep convective
dynamics set the pole-to-pole deformation-radius contrast.

%% ====================================================================
\section{Neptune and Uranus predictions}
\label{sec:neptune}

For Neptune ($R = 24{,}600$\,km,
$\Omega = 1.08 \times 10^{-4}$\,rad\,s$^{-1}$,
$\varepsilon = 0.017$), the packing bound dominates.
The deformation radius $R_d$ is poorly constrained
($1{,}500$--$2{,}500$\,km from thermal wind estimates);
the predicted $N$ depends sensitively on it:
\begin{center}
\small
\begin{tabular}{lccc}
\toprule
$R_d$ (km) & $r_{\mathrm{excl}}/R_{\mathrm{ring}}$ &
  $N_{\mathrm{pack}}$ & $N_{\mathrm{pred}}$ \\
\midrule
$1{,}500$ & $0.39$ & $7$ & $\le 7$ \\
$2{,}000$ & $0.52$ & $5$ & $5$ \\
$2{,}500$ & $0.65$ & $4$ & $4$ \\
\bottomrule
\end{tabular}
\end{center}
\noindent
using $R_{\mathrm{ring}} \approx 5{,}000$\,km (estimated from
the polar vortex extent) and $\beta$-drift factor $1.3$.
The prediction is $N \in \{3, 4, 5\}$---packing-dominated,
like Jupiter's south pole.
Neptune's oblateness ($\varepsilon = 0.017$) is negligible
compared to Saturn's ($\varepsilon = 0.098$); it shifts
$\lambda_{\min}$ by $< 0.002$ and does not affect $N_{\mathrm{crit}}$.

\smallskip
\noindent
For Uranus (axial tilt $98^\circ$): the extreme tilt
disrupts the Onsager condensation mechanism by preventing a
persistent polar jet.
The prediction is $N = 0$ (no polar polygon).

\smallskip
\noindent
\emph{Falsifiability}: a Neptune observation of $N \ge 6$
would violate the packing bound for any $R_d > 1{,}500$\,km;
a Uranus polar polygon would challenge the tilt-disruption
hypothesis.

%% ====================================================================
\section{Why the hexagon is uniquely two-dimensional}
\label{sec:uniquely-2d}

The hexagon ($N = 6$) is marginally stable in $d = 2$
($\lambda_{\min} = +0.5$) but unstable for $d \ge 3$
($\lambda_{\min} = -0.018$ at $d = 3$).
This follows from the generalised Casimir
$h(m, N, \Delta)$ of Paper~III,
Proposition~\ref{III-prop:generalised-casimir}:
\begin{center}
\small
\begin{tabular}{lcccl}
\toprule
$d$ & $\Delta = (d{-}2)/2$ & $N_{\mathrm{crit}}$ &
  $\lambda_{\min}(N{=}6)$ & Hexagon \\
\midrule
$2$ & $0$ & $7$ & $+0.50$ & Stable \\
$3$ & $1/2$ & $5$ & $-0.02$ & Unstable \\
$4$ & $1$ & $5$ & $-1.08$ & Unstable \\
$5$ & $3/2$ & $5$ & $-2.93$ & Unstable \\
\bottomrule
\end{tabular}
\end{center}
\noindent
Saturn's hexagon requires the logarithmic interaction
($d = 2$, $\Delta = 0$), which is the Green's function of
the 2D Laplacian.  In 3D, the Coulomb interaction ($1/r$)
destabilises the hexagon.
The logarithmic interaction survives in a quasi-2D
atmosphere because the weather layer is thin
($H_{\mathrm{eff}} \ll R_{\mathrm{planet}}$) and the
barotropic mode dominates: the effective dynamics is
2D to leading order, with 3D corrections entering at
$O(H_{\mathrm{eff}}/R)^2$.
This explains why hexagonal features are observed only on
gas giants with thin weather layers (Saturn) and not on
terrestrial planets or in 3D laboratory flows.

%% ====================================================================
\section{Broader observational context and predictions}
\label{sec:observations}
\label{sec:planets}

\subsection{Polygon quality metrics and cross-planetary comparison}

The two planetary polygon types require different quality metrics.
For \emph{discrete vortex rings} (Jupiter), the natural measure is the
geometric deviation from a perfect $N$-gon:
\begin{equation}\label{eq:sigma-geom}
  \sigma_{\mathrm{geom}} = \frac{1}{N}\sum_{k=1}^{N}
  \bigl|\!\ln|z_k / z_k^{\mathrm{ideal}}|\bigr|,
\end{equation}
where $z_k^{\mathrm{ideal}}$ are the best-fit regular $N$-gon vertices.
This measures positional scatter of discrete cyclones; small
$\sigma_{\mathrm{geom}}$ means the ring closely approximates the
Thomson equilibrium.

For \emph{continuous jet meanders} (Saturn), the vortex positions $z_k$
are not defined.  The appropriate measure is the fractional radial
displacement of the jet center:
\begin{equation}\label{eq:eps-obs}
  \varepsilon_{\mathrm{obs}} = A_n / \bar{R},
\end{equation}
where $A_n$ is the half-amplitude of the $n$-th Fourier mode of the
jet-center radial position and $\bar{R}$ is the mean ring radius.
For Saturn's hexagon ($n = 6$, $\bar{R} \approx 13{,}200\;\mathrm{km}$),
the observed $\pm 2^\circ$ latitude oscillation
\citep{Godfrey1988, Antunano2015} gives
$A_6 \approx 1{,}900\;\mathrm{km}$ and
$\varepsilon_{\mathrm{obs}} \approx 0.14$ (range $0.11$--$0.18$ across
measurement epochs).

These two metrics are \emph{not interchangeable}: applying
$\sigma_{\mathrm{geom}}$ to Saturn (by treating the meander peaks as
point positions) yields $0.50$, an order of magnitude larger than
Jupiter's $0.05$--$0.06$, but this comparison is physically
meaningless---it conflates meander amplitude with positional scatter.
Table~\ref{tab:planets} reports both metrics, using each where it is
appropriate.

\begin{table}[htbp]
\centering
\caption{Planetary polygon configurations.
$\kappa_0/\kappa$ is estimated from Juno imagery \citep{Adriani2018}
using cyclone area as a proxy ($\sim\!30\%$ systematic uncertainty).
``---'' indicates the quantity is not applicable.}
\label{tab:planets}
\begin{tabular}{lcccl}
\toprule
System & $N$ & Quality metric &
  $\kappa_0/\kappa$ & Selection mechanism \\
\midrule
Saturn hexagon    & 6 & $\varepsilon_{\mathrm{obs}} = 0.14$ & N/A (jet) & Rossby stationarity \\
Jupiter north     & 8 & $\sigma_{\mathrm{geom}} = 0.06$ & $\sim\!1$ & Thomson + Onsager \\
Jupiter south     & 5 & $\sigma_{\mathrm{geom}} = 0.05$ & --- & Packing (Prop.~\ref{prop:complementary}, item 3) \\
\bottomrule
\end{tabular}
\end{table}

\subsection{Broader observational context}

Beyond the three planetary configurations, the theoretical stability
boundary $\kappa_{\mathrm{crit}}(N)$ from
Proposition~Paper~II, Proposition~9.1 can be compared with laboratory and
numerical observations:

\emph{Laboratory experiments.}  \citet{Jansson2006} and
\citet{BarbosaAguiar2010} observe polygons with $N = 2$ through $6$
in rotating-tank experiments by varying the differential rotation rate.
All observed $N$ values fall in the $\kappa_{\mathrm{crit}} = 0$ range
($N \le 7$), consistent with the Thomson stability boundary.  No
laboratory experiment has produced a stable polygon with $N \ge 8$
without a central vortex---also consistent.  (The laboratory mechanism
is barotropic instability of a shear layer, not Thomson point vortex
dynamics, so the agreement is structural rather than quantitative.)

\emph{Numerical simulations.}  \citet{Siegelman2022} produce $N = 8$
(north) and $N = 5$ (south) vortex crystals from 2D turbulence on polar
caps, with a central polar cyclone emerging alongside the ring in both
cases.  The central cyclone provides the stabilization required by
Proposition~Paper~II, Proposition~9.1 for $N = 8$.
The \citet{Siegelman2022} simulations include continuous small-scale
forcing (representing convective injection) and large-scale dissipation,
yet the polygon crystal is the persistent stationary state: the system
is not frozen into a one-time equilibrium but continuously re-attracted
to it.  Within the 2D barotropic framework, this is numerical evidence
that the constrained energy minimum acts as a \emph{dynamical attractor}
under turbulent forcing.  The caveat (Section~\ref{sec:extensions}) is
that these are idealized 2D simulations, not a faithful model of
Jupiter's three-dimensional polar dynamics; the agreement supports the
Onsager picture but does not constitute a proof of its applicability.

\emph{Uranus.}  \citet{Akins2023} detect a single polar cyclone at
Uranus's north pole---a single-vortex ($N = 1$) configuration, below
the polygon threshold.  No ring structure has been observed.  Whether
the absence of a Uranian polygon reflects insufficient vortex formation
(macro-scale constraint) or insufficient observations (the pole has been
poorly resolved) remains open.

\emph{Neptune.}  No polar polygon has been observed, and Voyager~2
imagery is insufficient to resolve ring structure.  However, the
framework makes a falsifiable prediction.

\textit{Regime:}
At $84^\circ$N, $\beta_N = 2\Omega_N\cos(84^\circ)/R_N \approx
9.2\times10^{-13}$\,m$^{-1}$s$^{-1}$ and $R_{\mathrm{ring}} \approx
2{,}570$\,km, giving $\beta R^3 \approx 1.6\times10^7$\,m$^2$\,s$^{-1}$.
For Jovian-scaled cyclones (core radius $r_c \approx 240$\,km,
$V_{\max}\approx 80$--100\,m\,s$^{-1}$), $\kappa \approx
1.2\text{--}1.5\times10^8$\,m$^2$\,s$^{-1}$ and
\[
  \frac{\beta R^3}{\kappa}\approx 0.10\text{--}0.13 \ll 1.
\]
Neptune is firmly vortex-dominated---more so than Jupiter ($\beta
R^3/\kappa \approx 0.6$), because its smaller radius ($R_N \approx
R_J/3$) reduces $R_{\mathrm{ring}}^3$ by a factor of $\sim\!25$
despite a slightly larger $\beta$.  The Rossby stationarity mechanism
does not apply; the relevant bridge is Thomson/Onsager.

\textit{Prediction:}
The first-principles chain (Section~\ref{sec:first-principles})
narrows the Neptune prediction substantially.  Neptune's polar
ring radius is $R_{\mathrm{ring}} \approx 3{,}000$\,km
($R_N \sin 7^\circ$), which is $2.9\times$ smaller than Jupiter's.
Its Coriolis parameter at~$83^\circ$ is $f_N = 2.15\times
10^{-4}\;\mathrm{s}^{-1}$, only $0.62\times$ Jupiter's.
Consequently, the deformation radius
$R_d = \sqrt{g' H_{\mathrm{eff}}}/f$ is $1.6\times$ larger than
Jupiter's for the same atmospheric structure, and the packing ratio
$\varepsilon/R = R_d/R_{\mathrm{ring}}$ is $4.7\times$ tighter:
\begin{center}
\begin{tabular}{lcccc}
\toprule
$R_d$ & $\varepsilon/R$ & $N_{\mathrm{pack}}$ &
$N$ (no CV) & $N$ (with CV) \\
\midrule
$500$\,km (shallow) & $0.17$ & $18$ & $7$ & $8$ \\
$1{,}000$\,km (moderate) & $0.33$ & $9$ & $7$ & $8$ \\
$1{,}300$\,km & $0.43$ & $7$ & $7$ & $7$ \\
$1{,}500$\,km (Jupiter-like) & $0.50$ & $5$ & $5$ & $5$ \\
\bottomrule
\end{tabular}
\end{center}
Neptune's strong internal heat flux ($F_{\mathrm{int}} = 0.43\;
\mathrm{W\,m}^{-2}$, $2.6\times$ absorbed solar) suggests vigorous
convection, favouring $R_d \gtrsim 1{,}000$\,km.  In this regime
\textbf{packing is the binding constraint}, and the prediction
narrows to $N \in \{3, 4, 5\}$---a triangle, square, or pentagon.

The physical reason is geometric: Neptune's ring radius is too small
for many cyclones of Jupiter-like size to fit.  The $4.7\times$
tighter packing ratio (from the combination of smaller planet and
weaker Coriolis) means that even moderate-size cyclones crowd the
ring.  This is a qualitative difference from Jupiter north, where
Thomson (not packing) is the binding constraint.

This prediction is testable by a Neptune polar orbiter (e.g.\
Trident-class mission): imaging the polar region at
$\sim\!100$\,km resolution would resolve individual cyclones and
determine $N$ directly.

All observations are consistent with the theoretical stability
boundary: $N \le 7$ is observed without a central vortex, $N = 8$
is observed only with one, and no $N \ge 9$ ring has been reported
in any system.


%% ====================================================================
\section{Discussion and limitations}
\label{sec:discussion-new}

\subsection{Why 2D and not 3D}

The argument depends critically on two-dimensionality.  The 2D Green's
function $G(r) = -\ln r/(2\pi)$ is the unique dilation-invariant
pairwise interaction with $h' < 0$
(Remark~Paper~I, Remark~2.2; Corollary~Paper~I, Corollary~2.4
lifts this to the $N$-body level).  In 3D, the Green's function
$G(r) = -1/(4\pi r)$ has $G'(r) = 1/(4\pi r^2) > 0$---the sign
rule is violated, the $N$-gon is not an energy extremum along the
spiral, and the stability analysis gives different results.  The
present mechanism---sign rule plus energy-Casimir stability---does
not operate in 3D; whether other mechanisms can produce persistent
polygonal vortex structures in three dimensions is a separate
question.


\subsection{What is novel}

\noindent\textbf{Attribution.}
The Havelock eigenvalue $\lambda_m = (N{-}1) - m(N{-}m)/2$ and the
stability boundary $N \le 7$ are classical \citep{Havelock1931}.
Theorem~Paper~II, Theorem~8.2 bridges this to the Onsager
literature, where the polygon is often described as an energy
maximum: it is a constrained minimum, compatible with but distinct
from the unconstrained Onsager picture (Corollary~Paper~II, Corollary~8.3).

\smallskip\noindent
What is new:
(1)~the constrained-minimum identification and its
Lagrange-multiplier derivation;
(2)~the analytic sign rule (Paper~I, Theorem~3.1);
(3)~the three-scale separation and two-timescale Onsager--Arnold
resolution;
(4)~the constraint-intersection framework
(Proposition~\ref{prop:complementary});
(5)~the RG classification of the logarithm as IR fixed point
(Proposition~Paper~I, Proposition~2.1);
(6)~the spherical and hyperbolic stability boundaries
with exact thresholds;
(7)~the curvature-stability dichotomy and spectral flow
(Propositions~Paper~I, Proposition~5.5--Paper~I, Proposition~5.6);
(8)~the algebraic-integer classification $N \in \{8,9,11,15,23\}$,
golden ratio terminal case, binary subdivision, Pell--conformal
inversion equivalence, and continued fraction structure
(Section~Paper~I, \S6);
(9)~the blob bridge, $K_0$ central-vortex computation, and
circulation disorder resilience;
(10)~the first-principles prediction chain reducing planetary
$N$ to atmospheric parameters.

\smallskip\noindent\textbf{Geometry selects placement.}
A recurring theme of the series is that stability depends not
only on $N$ and the curvature, but on \emph{where} the polygon
sits.  On Saturn, the oblateness gradient makes the pole more
stable than lower latitudes: the hexagon exists at
$78^\circ$N but not below the critical latitude
$72^\circ$N (Section~\ref{sec:oblateness}), because the
Gaussian curvature decreases toward the pole and the binding
eigenvalue $\lambda_3$ flips sign at $\varphi_{\mathrm{crit}}$.
Paper~I, \S7 establishes the abstract counterpart:
on an arithmetic surface
$\Gamma\backslash\mathbf{H}^2$, the $L^2$-Morse index at
a palindromic threshold is \emph{fractional}---$61\%$ of the
Bolza surface is unstable while the automorphism-fixed centre
remains stable.
In both cases the mechanism is the same: a
\emph{symmetry-enhanced point} (the geographic pole on Saturn,
the Aut-fixed centre on the Bolza surface) concentrates the
stabilising curvature correction, making the vortex polygon
viable at that location when it would be unstable generically.
The oblateness of Saturn and the automorphism group of the
Bolza surface are playing the same mathematical role---they
select the placement.

\smallskip\noindent\textbf{Relation to prior work.}
\citet{BouchetVenaille2012} and \citet{ChavanisSommeria1996} classify
self-organized vortex states via statistical mechanics;
\citet{Aref2003} reviews vortex crystal stability.
Our contribution adds the energy-curvature sign rule, the
constrained-minimum clarification, and the planetary
constraint-intersection framework---none of which appear in
these references.  The $N = 7$ quartic coefficient $\alpha_0 = 45/14$
extends \citet{KurakinYudovich2002} (who proved $\alpha_0 > 0$
without the value) and complements \citet{Schmidt2004} (who computed
$9/28$ in Sylvester-normalized coordinates).


\subsection{Limitations and scope}

The mathematical results---the constrained energy minimum
(Paper~II, Theorem~8.2), the sign rule
(Paper~I, Theorem~3.1), and the stability table---are proven
for discrete point vortices.  They apply directly to Jupiter's polar
cyclone rings.  For Saturn's hexagon, which is a continuous jet meander,
the relevant mechanism is Rossby wave stationarity of the QGPV
equation, not the point vortex energy landscape.  The continuous and
discrete descriptions share the logarithmic Green's function; the
vortex blob analysis (Section~Paper~II, \S12) establishes that the
stability boundary is preserved under regularization and that the
$N = 7$ marginal case is resolved.

The Saturn bridge (Section~Paper~II, \S12.3) connects the
continuous QGPV eigenvalue problem to the point vortex Hessian via
critical-layer concentration.  The key mechanism is
$\mathbb{Z}_N$~symmetry: for any $\mathbb{Z}_N$-symmetric cat's-eye
structure, the lobe circulations are equal and the centroids lie on
the $N$-gon vertices, so the constrained Hessian has the same inertia
as the Havelock Hessian (Remark~Paper~II, Remark~12.3).  This is verified
numerically for both smooth (Stuart vortex) and step-PV
(Prandtl--Batchelor) cat's-eye models.  The braid vorticity fraction
vanishes as $O(1{-}\eta)$ in the Stuart concentration limit
(Remark~Paper~II, Remark~12.3), satisfying the conditions of
Proposition~Paper~II, Proposition~12.2.

The quantitative bridge at Saturn's actual parameters can be closed
as follows.  The relevant concentration parameter is the
\emph{physical} jet width divided by the hexagon radius:
$\varepsilon/R = \sigma_{\mathrm{jet}}/R_{\mathrm{hex}}
= 2.5 \times 10^6/(5.1 \times 10^7) = 0.049$,
not the log-polar width $\sigma_{\mathrm{lp}} \approx 0.15$
sometimes quoted.  With $d_{\min} = 2\sin(\pi/6) = 1$, the
concentration condition~Paper~II, eq.~12.5 requires
$\varepsilon/d_{\min} < \delta(6) = 0.096$.  The physical ratio
$0.049 < 0.096$ satisfies this condition.
The Weyl perturbation bound gives
$\|\nabla^2\!H_\varepsilon - \nabla^2\!H_0\|_{\mathrm{op}}
\le 54\varepsilon^2 = 0.13$, which is $26\%$ of the spectral gap
$\lambda_{\min}(6) = 1/2$---absorbed with a $3.9\times$ margin.
The actual $O(\varepsilon^2)$ blob correction is \emph{stabilizing}
($c_6 = 13/12 > 0$, shifting $\lambda_{\min}$ upward by $+0.003$),
and the separatrix vorticity sheet contributes $O(\varepsilon^4)$
to the eigenvalue (quadrupole coupling of a thin dipole layer),
negligible compared to the gap.  Saturn's finite jet width therefore
\emph{strengthens}, not weakens, the $N = 6$ stability.

The Saturn hexagon is therefore supported by three arguments:
(i)~Rossby stationarity (Proposition~\ref{prop:complementary}) selects
$n = 6$ independently of the energy landscape;
(ii)~geometric consistency with the spherical stability boundary
(Section~Paper~I, \S5.1); and
(iii)~the critical-layer bridge (Remark~Paper~II, Remark~12.3), which
establishes Hessian inertia preservation for $\mathbb{Z}_N$-symmetric
cat's-eye structures---now verified quantitatively at Saturn's
physical parameters ($\varepsilon/d_{\min} = 0.049 < \delta(6) = 0.096$,
Weyl perturbation $26\%$ of spectral gap).  All three arguments
converge on $N = 6$.

The Onsager negative-temperature hypothesis is a physical assumption
whose applicability to Jupiter's polar region is not directly tested
here.  It is the key step that distinguishes $N = 8$ from $N = 7$
in the selection argument (Proposition~\ref{prop:onsager-n-select}):
without negative temperature, the framework identifies $N = 8$ as
\emph{stable} but does not uniquely select it over $N = 7$.

The \citet{Siegelman2022} simulations provide the most direct evidence
available.  Their 2D turbulence on polar caps develops an inverse energy
cascade that drives kinetic energy to the largest scales consistent with
the polar geometry---the signature of a Onsager negative-temperature
state.  The polygon crystal is the maximum-$H$ configuration within
the Thomson-stable set: starting from random initial conditions, the
simulations converge to the polygon ring with a central cyclone, which
is the $N_{\max}$ configuration under Proposition~\ref{prop:onsager-n-select}
(they observe $N = 8$ north and $N = 5$ south, matching the Juno
observations; the south-pole discrepancy from the Thomson/Onsager
prediction is attributed to geometric packing as discussed above).
This is indirect but consistent evidence that the planetary system is
near negative temperature.  A direct test would require measuring the
energy spectrum at the polar scale in Juno or a dedicated mission.

Cassini and Juno data are digitized from published figures.

\emph{Dynamical maintenance.}%
\label{sec:dynamical-maintenance}
The constrained energy minimum (Theorem~Paper~II, Theorem~8.2) is
a static result: it identifies the polygon as the lowest-energy
arrangement on the conserved-quantity surface, and by Arnold's energy-Casimir
stability theorem \citep{Arnold1966} this minimum is Lyapunov stable
under the conservative 2D Euler dynamics---perturbations stay bounded.
Whether this translates into an attractor under convective forcing
depends on the ratio $\mathcal{R} = \tau_{\mathrm{inject}}/\tau_{\mathrm{restore}}$.

For Jupiter's $N=8$ ring (Juno parameters: $R_{\mathrm{ring}}\approx
7{,}500\;\mathrm{km}$, $\kappa_0/\kappa \approx 1$;
\citealt{Adriani2018}), we use the full-cyclone circulation
$\kappa = 2\pi r_c V_{\max} \approx 1.6\times10^9\;\mathrm{m^2\,s^{-1}}$
with $r_c\approx3{,}200$\,km (the full Juno-observed cyclone radius)
and $V_{\max}\approx80$\,m\,s$^{-1}$.  This is the appropriate measure
of total vortex strength for computing orbital periods and eigenfrequencies.
It differs from the Rankine-core estimate
$\kappa_{\mathrm{Rank}}\approx3.5\times10^8$\,m$^2$\,s$^{-1}$
(using only $r_c\approx700$\,km, the solid-body rotation region)
that enters the regime number $\beta R^3/\kappa$ in
Section~\ref{sec:extensions}: the regime criterion requires only a
characteristic vorticity scale, for which the core is sufficient.
\begin{itemize}
\item \emph{Restoring:} The ring orbital period is $T_{\mathrm{orb}} =
  8\pi^2 R^2/[\kappa(N{-}1)] \approx 4.6\;\mathrm{days}$.
  The most marginal mode ($m=4$) has Lagrangian eigenvalue
  $\lambda_{\mathrm{eff}} \approx 3$ (bare eigenvalue $-1$, raised to
  $\approx+3$ by the central cyclone with $\kappa_0/\kappa=1$).
  The oscillation e-folding time is
  $\tau_{\mathrm{restore}} = (\sqrt{\lambda_{\mathrm{eff}}}\,\Omega)^{-1}
  \approx 0.4\;\mathrm{days}$.

\item \emph{Injection:} The upper-bound convective power into one
  cyclone's footprint is $P = F_{\mathrm{int}}\,\pi r_c^2$, where
  $F_{\mathrm{int}} \approx 7.5\;\mathrm{W\,m^{-2}}$ is Jupiter's
  interior heat flux \citep{Guillot2005}; with $r_c = 700\;\mathrm{km}$,
  $P \approx 1.15\times10^{13}\;\mathrm{W}$.
  The tropospheric kinetic energy per cyclone is estimated as
  $E_K \approx \tfrac{1}{2}\rho V_{\max}^2\,\pi r_c^2 H$
  with $\rho \approx 0.2\;\mathrm{kg\,m^{-3}}$, $H \approx 40\;\mathrm{km}$,
  $V_{\max} \approx 80\;\mathrm{m\,s^{-1}}$:
  $E_K \approx 4\times 10^{19}\;\mathrm{J}$, giving
  $\tau_{\mathrm{inject}} = E_K/P \approx 40\;\mathrm{days}$.
\end{itemize}
The ratio $\mathcal{R} \approx 40/0.4 \approx 100$: the restoring force
operates $\sim\!100\times$ faster than the upper-bound convective
injection.  Including the inverse-cascade efficiency (typically
$\lesssim 1\%$ of the heat flux reaches cyclone scales) raises
$\mathcal{R}$ to order $10^4$.  The Juno four-year observational
baseline provides a direct empirical lower bound:
$\mathcal{R} \gtrsim 4\;\mathrm{yr}/0.4\;\mathrm{days} \approx 3500$.
The polygon is firmly in the restoring-dominated regime; convective
forcing is a persistent perturbation that the energy minimum absorbs,
not a mechanism that competes with the minimum for selection of the
configuration.

\emph{Two-timescale resolution of the Onsager--Arnold tension.}
The constrained energy minimum (Lyapunov stability) and the Onsager
negative-temperature selection (ergodic exploration of the energy
landscape) appear to be in tension: if the polygon is dynamically
stable, how does the system ergodically explore away from it?
The resolution is that they operate \emph{sequentially}, not
simultaneously:
\begin{itemize}
  \item \textbf{Formation phase}
    ($\tau_{\mathrm{form}} \sim$\,weeks to months):
    The inverse cascade concentrates same-sign vorticity into $N$
    coherent clusters.  During this phase the system is turbulent
    and ergodic; the Onsager negative-temperature argument determines
    \emph{which} $N$ is selected ($N_{\max}$, the highest-energy
    Thomson-stable configuration).  The polygon has not yet formed;
    the clusters are accumulating and arranging.
  \item \textbf{Persistence phase}
    ($\tau_{\mathrm{persist}} \sim$\,years to decades):
    Once the $N$ clusters reach the polygon arrangement, Arnold
    stability takes over.  The polygon is a constrained energy
    minimum with $\mathcal{R} \approx 100$; perturbations are
    bounded and the system is \emph{no longer ergodic}---it is
    trapped in the stability basin.  The polygon persists because
    it is dynamically stable, not because of ongoing statistical
    equilibrium.
\end{itemize}
The Onsager argument selects $N$ \emph{before} the polygon forms;
the Arnold argument maintains it \emph{after}.  The ratio
$\mathcal{R} \approx 100$ quantifies the transition: the dynamical
restoring timescale ($0.4$~days) is much shorter than the formation
timescale ($\sim\!40$~days), so once the polygon assembles, it
freezes into the energy minimum essentially irreversibly on
observational timescales.

\emph{Moist convection as vorticity source.}
The leading mechanism proposed for maintaining Jupiter's polar cyclones
in the observational literature \citep{ONeill2015} is moist convective
plumes injecting cyclonic vorticity preferentially at high latitudes.
The geometric argument is that near the pole, the $\beta$-effect
preferentially suppresses anticyclonic plumes (their vorticity anomaly
is opposed by the planetary background), while cyclonic plumes are
unimpeded; the result is a net positive-vorticity injection into the
polar cap.  This mechanism is quite different from---and more specific
than---the Onsager inverse-cascade picture invoked here.

The two mechanisms are complementary rather than competing, operating
at different scales and addressing different aspects of the problem:
\begin{itemize}
\item \textbf{O'Neill et al.\ generation}: moist convection injects
  cyclonic PV into the polar region, selecting \emph{sign} (cyclones,
  not anticyclones).  This is the vorticity \emph{source}.
\item \textbf{Onsager organization}: the injected cyclonic PV undergoes
  an inverse energy cascade, clustering into $N$ coherent vortices.
  This is scale-blind to the source mechanism.
\item \textbf{Thomson/energy landscape}: the $N$ clusters arrange into
  a regular polygon.  This is determined entirely by the point-vortex
  energy function.
\end{itemize}
The Onsager inverse cascade (and the present paper) requires only that
same-sign vorticity be present in the polar region; it does not predict
the sign.  The O'Neill mechanism provides that sign selection.  Without
it, the Onsager picture would predict equal probability of cyclonic and
anticyclonic polar clusters---which is not observed.  A complete theory
of Jupiter's polar polygons therefore requires both mechanisms:
O'Neill for the vorticity source and polarity, and the present
Thomson/Onsager framework for the organization and stability of
whatever is produced.

\subsection{Toward first-principles prediction}
\label{sec:first-principles}

The constraint-intersection framework
(Proposition~\ref{prop:complementary}) predicts $N$ given
observational inputs about the vortices themselves ($\kappa_0/\kappa$,
$r_{\mathrm{cyclone}}$, $R_{\mathrm{ring}}$).  A stronger result
would predict $N$ from \emph{atmospheric} parameters that are
independently measurable without observing the polygon structure.
The deformation-radius chain partially achieves this reduction:
\begin{enumerate}
  \item \emph{Planetary parameters} $(\Omega, R, \varphi)$ give
    $f = 2\Omega\sin\varphi$ and
    $\beta = 2\Omega\cos\varphi / R$ (exact).
  \item \emph{Atmospheric structure} $(g', H_{\mathrm{eff}})$---the
    reduced gravity and effective weather-layer depth, measurable from
    temperature profiles and gravity data---give the deformation
    radius $R_d = \sqrt{g' H_{\mathrm{eff}}}/f$ (derived), which
    sets the cyclone core size $r_{\mathrm{core}} \sim R_d$.
  \item \emph{Domain admissibility}
    $\sin(\pi/N) \ge R_d / R_{\mathrm{ring}}$ gives
    $N_{\mathrm{pack}}$ (derived from the Laplacian).
  \item \emph{Thomson + Onsager} gives $N_{\mathrm{Thomson}} = 8$
    for $\kappa_0/\kappa \ge \tfrac{1}{2}$, which is the
    generic outcome of Onsager condensation when a central vortex
    is present (derived).
  \item $N = \min(N_{\mathrm{pack}},\, N_{\mathrm{Thomson}})$
    (constraint intersection).
\end{enumerate}
For Saturn, the Rossby path replaces steps~2--4:
$n^* = R_{\mathrm{hex}}\sqrt{\beta/U_{\max}}$, where
$U_{\max}$ is the jet speed (atmospheric observation, not a
vortex observation).

At Jupiter, the north--south asymmetry ($N = 8$ vs.\ $N = 5$) reduces
to a single parameter: $\chi = R_d^{(\mathrm{south})}/R_d^{(\mathrm{north})}$,
the ratio of deformation radii at the two poles.  With
$R_d^{(\mathrm{north})} \approx 2{,}500$\,km (matching the observed
north-pole cyclone radius at $H_{\mathrm{eff}} \approx 3{,}000$\,km,
$g' \approx 0.3$\,m\,s$^{-2}$), the constraint intersection gives
$N = 8$ (Thomson-limited).  At the south pole,
$\chi \approx 1.34$ gives $R_d^{(\mathrm{south})} \approx 3{,}350$\,km;
with the $\beta$-drift exclusion factor $\approx 1.3$
\citep{Gavriel2021}, the effective exclusion radius
$r_{\mathrm{excl}} \approx 4{,}350$\,km yields
$N_{\mathrm{pack}} = 5$ (packing-limited).  The required
$\chi \approx 1.34$ matches the observed cyclone size ratio
$r_{\mathrm{cyclone}}^{(\mathrm{south})}/r_{\mathrm{cyclone}}^{(\mathrm{north})}
\approx 3{,}500/2{,}500 = 1.40$ to within $5\%$.

What remains is to derive $\chi$ itself from atmospheric physics.
The south pole receives $\approx 22\%$ more insolation at perihelion
(Jupiter's orbital eccentricity $e = 0.049$), but radiative scaling
gives only $\chi \sim 1.07$---insufficient.  A complete derivation
of $\chi$ from moist convective dynamics remains an open problem;
the framework reduces the prediction to this single parameter.

Three configurations from two planets is consistent with the proposed
mechanism but not sufficient to establish universality.  Uranus and
Neptune observations would substantially strengthen the case.  We use
``consistent with a general mechanism'' rather than ``universal'' to
reflect the available evidence.

\subsection{Falsifiability}

The framework makes five testable predictions whose violation would
require fundamental revision:
\begin{enumerate}
\item A \emph{stable $N \ge 9$ polar polygon} on any gas giant
  would violate the Thomson bound
  ($\kappa_{\mathrm{crit}}(9) = 1$, requiring the central vortex
  to exceed the ring vortex---an extreme configuration not produced
  by Onsager condensation).
\item \emph{Saturn's hexagon below $72^\circ$N latitude}
  would violate the oblateness sign-flip
  ($\lambda_3 < 0$ below $\varphi_{\mathrm{crit}}$).
\item A \emph{hexagonal feature in a 3D-dominated regime}
  ($N_{\mathrm{crit}}$ drops from $7$ to $5$ at $d = 3$;
  the hexagon is uniquely two-dimensional).
\item \emph{Neptune $N \ge 6$} would violate the packing bound
  for any deformation radius $R_d > 1{,}500$\,km.
\item \emph{Jupiter $N = 8$ in shallow-layer models}
  ($R/R_d > 1.5$) would violate the $K_0$ suppression that
  requires a deep weather layer.
\end{enumerate}

%% ====================================================================
\section{Stellar applications}
\label{sec:stellar}

The polygon stability framework extends beyond gas giants to any
system with two-dimensional vortex dynamics and a logarithmic
(or logarithmic-like) interaction.

\subsection{Neutron star superfluid vortices}

The interior of a neutron star contains a superfluid of neutron
Cooper pairs.  Rotation at angular frequency~$\Omega$ creates
quantised vortices with circulation
$\kappa = h/(2m_n) \approx 2.0 \times 10^{-7}$\,m$^2$\,s$^{-1}$
and areal density $n_v = 2\Omega/\kappa$.

\begin{center}
\small
\begin{tabular}{lccc}
\toprule
Pulsar & $f$ (Hz) & Vortex density (m$^{-2}$) & Spacing ($\mu$m) \\
\midrule
Crab (young) & $30$ & $1.9 \times 10^9$ & $23$ \\
Vela (glitching) & $11$ & $7.0 \times 10^8$ & $38$ \\
MSP J1748 & $716$ & $4.6 \times 10^{10}$ & $4.7$ \\
\bottomrule
\end{tabular}
\end{center}

\noindent
The superfluid vortices interact via the two-dimensional
logarithmic Green's function---precisely the Havelock kernel.
Three predictions follow directly:

\begin{enumerate}
\item \emph{Triangular vortex lattice.}
  Theorem~\ref{III-thm:triangle-universality} (triangle universality)
  predicts that the vortex array forms a triangular lattice---the
  unique stable regular tiling for the logarithmic interaction.
  This is consistent with theoretical expectations from the
  Abrikosov lattice in type-II superconductors.

\item \emph{Ring instability at $N = 7$.}
  Any ring of $N > 7$ co-rotating vortices at the crust--core
  boundary is unstable by the Havelock theorem and will fragment.
  Rings of $N \le 7$ can persist as metastable configurations.

\item \emph{Quantised glitch sizes.}
  Pulsar glitches (sudden spin-up events) arise from the
  collective unpinning of vortices at the crust--core boundary.
  If the unpinning occurs in ring-shaped clusters, the
  Havelock stability limit constrains the cluster size to $N \le 7$.
  The angular momentum transfer per ring:
  \[
    \Delta\Omega_N = \frac{N\kappa}{2\pi R_{\mathrm{pin}}^2}
  \]
  is \emph{quantised} in units of $N = 3, 4, 5, 6, 7$.
  Glitch magnitudes should cluster around these values,
  with a maximum single-ring contribution at $N = 7$.
  The $7 \to 8$ instability means that larger clusters
  break into sub-rings before unpinning.
\end{enumerate}

\subsection{The solar supergranulation}

The Sun's surface exhibits ${\sim}150$ supergranulation cells
at the equator (cell size ${\sim}30{,}000$\,km).
With $N \approx 150 \gg 7$: the equatorial ring of cells
is \emph{deeply unstable} by the Havelock criterion.
This is consistent with the observed turbulent, time-varying
character of the supergranulation pattern---the Sun has no
stable polygon because $N \gg N_{\mathrm{crit}}$.

\subsection{Rapidly rotating stars}

Stars with extreme oblateness ($\varepsilon > 0.5$, such as
Vega, Altair, and Achernar) have surface curvature gradients
analogous to Saturn's oblateness.  If convective or magnetic
vortices form at the poles of these stars, the oblateness
correction shifts $N_{\mathrm{crit}}$ (as for Saturn's hexagon,
\S\ref{sec:oblateness}).  For $\varepsilon > 0.5$:
$N_{\mathrm{crit}}$ may increase by $1$--$2$ relative to a
sphere, potentially stabilising polar polygons that would be
unstable on a non-rotating star.
This is a prediction for future direct-imaging observations
of rapid rotators with extreme-adaptive-optics instruments.

\subsection{Magnetic white dwarfs}

In magnetic white dwarfs ($B \sim 10^6$--$10^9$\,G), the
electron gas supports quantised magnetic flux tubes
($\Phi_0 = h/(2e) = 2.07 \times 10^{-15}$\,Wb).
When the London penetration depth exceeds the inter-tube
spacing, the flux tubes interact via the two-dimensional
logarithmic kernel---the Havelock regime.
The tessellation stability theorem predicts a triangular
flux tube lattice, and ring configurations of flux tubes
are limited to $N \le 7$ by the Havelock stability boundary.

\bibliographystyle{plainnat}
\bibliography{../shared/refs}
\end{document}
